\subsection{A proved family of nonabelian chief lifts}

\begin{statement}{D.28}{Proved}{Generating relative lifts for transitive $A_5$ wreath kernels}
Let $D\leq\operatorname{Sym}(\Omega)$ be a transitive finite
$5$-group, let $n=|\Omega|\geq25$, and let
$q:\widehat F=\langle x_1,x_2\rangle\twoheadrightarrow D$
be an epimorphism.  Put $x_3=(x_1x_2)^{-1}$ and $d_i=q(x_i)$.
Assume that each of the three permutations $d_i$ has no fixed
point on $\Omega$.  Give $\Omega$ its right $D$-action, and
write
\[
 M=A_5^{\Omega},\qquad
 (d\cdot a)(\omega)=a(\omega d),\qquad W=M\rtimes D.
 \tag{D.28.1}
\]
For every three elements $U_i\in W$ above $d_i$ there exist
$a_i\in M$ such that
\[
 V_i=a_iU_ia_i^{-1},\qquad
 V_1V_2V_3=1,\qquad\langle V_1,V_2\rangle=W.
 \tag{D.28.2}
\]
Equivalently, every such finite relative lifting problem has
a generating solution, with the three base conjugators fixed.
The kernel $M$ is a nonabelian chief factor: it is a product
of nonabelian simple groups, and $D$ permutes its factors
transitively.

In particular this applies to the regular action of every
noncyclic two-generated finite $5$-group $D$.  It also has
the following exact coupled consequence.  Let $L$ be a finite
$5$-group with epimorphism $\eta:L\to D$, and suppose
specified epimorphisms from $\widehat F$ to $W$ and $L$
have the same projection to $D$.  Let
$\beta\in\operatorname{Aut}(L)$ induce an automorphism
$\alpha\in\operatorname{Aut}(D)$, and suppose its three
source peripheral images satisfy
\[
       \beta(l_i)=h_i l_i^\lambda h_i^{-1}
       \quad(1\leq i\leq3),\qquad
       \lambda\in\widehat\Z^\times.
 \tag{D.28.3}
\]
Then $\beta$ has a generating peripheral lift from
$\widehat F$ to the actual fibre product $W\times_D L$,
over the map $\beta q_L$, with the same multiplier
$\lambda$ and with the given $h_i$ retained in the
$L$-coordinates of its conjugators.
\end{statement}

\begin{proof}
We first record all finite-group input used in the proof.
Exact enumeration of the sixty even permutations of five
letters gives the following commutator counts:
\[
\begin{array}{c|rrrrr}
 \text{element order}&1&3&2&5&5\\
 \text{conjugacy-class size}&1&20&15&12&12\\
 \#\{(s,t):[s,t]=u\}&300&63&32&65&65.
\end{array}
 \tag{D.28.4}
\]
In particular every element of $S=A_5$ is a commutator.
The complete subgroup enumeration has fifty-nine subgroups;
the only normal subgroup orders are $1,60$.  Its maximal
proper subgroups consist of five of order twelve, six of
order ten, and ten of order six.  Thus $S$ is nonabelian
simple and perfect, and every proper subgroup is contained
in one of twenty-one subgroups of size at most twelve.
Enumeration of every possible image of a fixed generating
pair gives $|\operatorname{Aut}(S)|=120$; the sixty inner
automorphisms are distinct, so $Z(S)=1$ and
$\operatorname{Out}(S)\simeq C_2$.  These are finite
certificates, described at the end of the proof.

Choose $\omega_0\in\Omega$, and put
\[
 H=q^{-1}(D_{\omega_0}),\qquad
 K=\ker q=\operatorname{core}_{\widehat F}(H).
 \tag{D.28.5}
\]
The equality with the core follows from the faithful
permutation action $D\leq\operatorname{Sym}(\Omega)$.
The subgroup $H$ has index $n$ and is the profinite
completion of the corresponding degree-$n$ cover of the
oriented three-boundary sphere, as in D.26.
Let its genus be $g$ and the number of its boundary
components be $b$.  Every cycle of any $d_i$ has length
at least five, so
\[
 b\leq\frac{3n}{5},\qquad
 2g=n+2-b,\qquad
 2g-2\geq\frac{2n}{5},\qquad g\geq6.
 \tag{D.28.6}
\]

We prove a uniform representation statement for this cover.
Prescribe an arbitrary conjugacy class $C_j$ of $S$ for
each of its $b$ oriented boundary loops.  Put
$P_C=\prod_j|C_j|\geq1$.
The surface presentation
\[
 \left\langle A_1,B_1,\ldots,A_g,B_g,c_1,\ldots,c_b
  \ \middle|\ \prod_{r=1}^{g}[A_r,B_r]\prod_{j=1}^{b}c_j=1
 \right\rangle
 \tag{D.28.7}
\]
shows that the number of continuous homomorphisms
$\varphi:H\to S$ taking $c_j$ into $C_j$ is at least
\[
                         60^{2g-2}P_C.
 \tag{D.28.8}
\]
Indeed choose all $b$ boundary values and the first $g-1$
handle pairs arbitrarily.  The required value of the last
commutator has a solution by \textup{(D.28.4)}.  Each such
choice gives at least one distinct homomorphism; all maps
are continuous because their target is finite and $H$ is
the completion of the displayed finitely generated group.

The maps with proper image in $S$ number at most
\[
                         21\cdot12^{2g}P_C.
 \tag{D.28.9}
\]
For each maximal proper subgroup, choose its handle values
and the boundary values in its intersections with the
$C_j$, and ignore the product relation to obtain an upper
bound.  Summing over the twenty-one maximal subgroups
gives \textup{(D.28.9)}.

Among the surjective maps we must also exclude repeated
coordinate kernels.  For such a $\varphi$, the restriction
$\varphi|_K$ is onto $S$: the quotient of $S$ by its
normal subgroup $\varphi(K)$ is a quotient of the
$5$-group $H/K$, whereas $S$ is perfect.
Put $N=\ker(\varphi|_K)$ and
$J=N_{\widehat F}(N)$.  One has $H\leq J$.
If $J>H$, conjugation on $K/N\simeq S$ defines a
continuous map $J\to\operatorname{Aut}(S)$.
Its composite to $\operatorname{Out}(S)$ factors through
the $5$-group $J/K$, and is consequently trivial.
Centerlessness therefore gives an actual homomorphism
\[
 \widetilde\varphi:J\longrightarrow S,
 \qquad
 \operatorname{Inn}(\widetilde\varphi(t))
       =\bigl(kN\longmapsto tkt^{-1}N\bigr).
 \tag{D.28.10}
\]
On $H$ this equals $\varphi$, since their inner
automorphisms of the surjective quotient $K/N$ agree.
It is the unique extension: an image of $t$ in any
extension is determined by its conjugation on $S$,
whose center is trivial.

There is an intermediate open subgroup
$H<H'\leq J$ with $[H':H]=5$.  To justify this step,
apply the normalizer property of finite $5$-groups to
$H/K<J/K$: a proper subgroup is strictly contained in
its normalizer.  That property follows by induction from
the center: either a central element is outside the
subgroup, or one passes to the quotient by the center.
Choose a subgroup of order five in the resulting
nontrivial normalizer quotient and take its inverse
image.  This gives the required $H'$.
It is determined by the set of its five right $H$-cosets,
which contains the identity coset.  There are at most
$\binom{n-1}{4}$ possibilities for $H'$.
Each is free profinite of rank $1+n/5$, and hence has
exactly $60^{1+n/5}$ homomorphisms to $S$.
Every surjective $\varphi$ with $J>H$ is the restriction
of one of these homomorphisms.  The number of such maps
is therefore at most
\[
                    \binom{n-1}{4}60^{1+n/5}.
 \tag{D.28.11}
\]
For a regular action one may replace the binomial factor
by the number of order-five subgroups of $D$, at most
$(n-1)/4$; the proof below uses the larger bound so it
also covers the stated transitive actions.

The two bad-map bounds are strictly smaller than
\textup{(D.28.8)}.  After division by that lower bound,
they are at most
\[
 \frac{75600}{5^{2g}}+
        \binom{n-1}{4}60^{1-n/5}
 \leq\frac{75600}{5^{12}}+
        \frac{\binom{24}{4}}{60^4}
 =\frac{38122819}{33750000000}<\frac1{500}<1.
 \tag{D.28.12}
\]
For the second inequality, the binomial expression
decreases when $n$ is increased by five from any
$n\geq25$: its successive ratio is bounded by
$(29/21)^4/60=707281/11668860<1$.
All possible $n$ are powers of five, so this estimate
applies to them.  We have proved that every boundary
class pattern has a surjective $\varphi$ with
$N_{\widehat F}(\ker\varphi|_K)=H$.

We now give the exact nonabelian Shapiro map from this
surface representation to the required source lift.
Choose $t_\omega\in\widehat F$ representing the right
coset $H t_\omega$, with $t_{\omega_0}=1$, so that
$\omega_0 q(t_\omega)=\omega$.  Define
\[
 c_\varphi(g)(\omega)=
       \varphi(t_\omega g t_{\omega q(g)}^{-1}),
 \qquad
       f_\varphi(g)=(c_\varphi(g),q(g))\in W.
 \tag{D.28.13}
\]
The element inside $\varphi$ belongs to $H$.
Factoring $t_\omega gh t_{\omega q(gh)}^{-1}$ at
$t_{\omega q(g)}$ gives exactly
\[
 c_\varphi(gh)(\omega)=
 c_\varphi(g)(\omega)c_\varphi(h)(\omega q(g)),
 \tag{D.28.14}
\]
so $f_\varphi$ is a continuous homomorphism over $q$.
For a cycle of $d_i$ of length $e$, starting at
$\omega$, its cycle product is
\[
 \prod_{r=0}^{e-1}c_\varphi(x_i)(\omega d_i^r)
       =\varphi(t_\omega x_i^e t_\omega^{-1}).
 \tag{D.28.15}
\]
The elements on the right are exactly the induced
oriented boundary generators in \textup{(D.28.7)},
up to their specified change of base paths.

For completeness, conjugation of $(v,d_i)$ by
$a\in M$ changes its function to
$a(\omega)v(\omega)a(\omega d_i)^{-1}$.
On each cycle its product is conjugated by the value
$a(\omega)$ at the chosen starting point.  Conversely,
a conjugator between the two cycle products determines
all its other values on that cycle by the recursion
\[
 a(\omega d_i)=v'(\omega)^{-1}a(\omega)v(\omega).
 \tag{D.28.16}
\]
The final equation is precisely the cycle-product
conjugacy.  Thus $M$-conjugacy classes above $d_i$ are
exactly the independently labelled $S$-conjugacy classes
of these cycle products.  Prescribe in
\textup{(D.28.7)} the classes arising from $U_i$.
Equations \textup{(D.28.15)--(D.28.16)} give
$f_\varphi(x_i)=a_iU_ia_i^{-1}$, with actual
$a_i\in M$, proving the required peripheral conditions.

It remains to check generation, not just existence of
a weak lift.  For $k\in K$, the coordinate maps of
$c_\varphi(k)$ are
$k\mapsto\varphi(t_\omega k t_\omega^{-1})$.
They are onto $S$, with kernels
$t_\omega^{-1}Nt_\omega$.  Two such kernels coincide
exactly when $t_{\omega'}t_\omega^{-1}\in J$.
Since $J=H$, this happens exactly for the same right
$H$-coset.  The kernels are therefore pairwise distinct.
The subdirect-product proof in D.25 now gives
$f_\varphi(K)=M$: a missing simple factor would force
one of the component epimorphisms to be an automorphic
copy of another and hence give equal kernels.
Together with $q(\widehat F)=D$, this proves
$f_\varphi(\widehat F)=W$, completing
\textup{(D.28.2)}.

If $D$ is noncyclic and its action is regular, each
$d_i$ is nonidentity: $d_1=1$, $d_2=1$, or
$d_3=1$ would make the two-generated group cyclic.
Every cycle then has the nontrivial $5$-power length
$|d_i|\geq5$, and $n=|D|\geq25$.
This proves the regular-case assertion.

We finally retain the prescribed base conjugators in
the coupled application.  Let $X_i\in W$ and
$l_i\in L$ be the source peripheral triples.  Choose
$B_i\in W$ above $\eta(h_i)$ and put
$U_i=B_iX_i^\lambda B_i^{-1}$.
These elements lie above $\alpha(d_i)$.
The three new permutations have the same cycle lengths
as the original $d_i$: they are conjugate to their
unit powers.  Apply \textup{(D.28.2)} to the source
map $\alpha q$ to obtain an epimorphism
$f_W:\widehat F\to W$ with
$f_W(x_i)=a_iU_ia_i^{-1}$.
The map
\[
 \widehat F\longrightarrow W\times_D L,
 \qquad g\longmapsto(f_W(g),\beta q_L(g))
 \tag{D.28.17}
\]
is defined because both components project to
$\alpha q(g)$.  Its peripheral conjugators are
the literal pairs $(a_iB_i,h_i)$, whose images in
$D$ agree, so it retains all three $h_i$.

To prove surjectivity of \textup{(D.28.17)}, write
$K_L=\ker q_L$ and $K_D=\ker q$.
The image of $K_D$ in $W$ under $f_W$ is $M$.
Its quotient by the image of $K_L$ is a quotient of
$K_D/K_L\simeq\ker\eta$, a $5$-group, and also a
quotient of the perfect group $M$; it is trivial.
Thus \textup{(D.28.17)} contains the whole kernel
$M$ over $L$, and maps onto $L$, which proves
surjectivity.  The same argument proves that the
original paired source is an epimorphism.  In these
fibre products the projection to $L$ is the maximal
pro-five quotient, since every map to a finite
$5$-group kills the perfect kernel $M$.

The program \texttt{a\_peripheral\_a5\_wreath.py}
provides all finite constants in this proof.  It
enumerates all $3600$ commutators and all subgroups by
successively adjoining a group element, yielding the
complete subgroup lattice and its normal and maximal
subgroups.  It tests all $60^2$ possible images of the
fixed generating pair, propagating the map along all
generator edges, and saves all $120$ automorphisms.
It checks \textup{(D.28.12)} and its monotonic ratio
by exact fractions.  For the first regular example,
$D=C_5^2$ with source directions $(1,0),(0,1),(-1,-1)$,
it records all three permutations of its twenty-five
points, their five cycles each, the genus six, and
all six order-five subgroups.  The source, output and
JSON certificate are retained; the output ends with
\texttt{PASS\_UNIFORM\_A5\_WREATH\_LIFTING\_CONSTANTS}.
\end{proof}

The no-fixed-point and degree hypotheses in D.28 are
used in the exact genus bound \textup{(D.28.6)}.
The admitted family includes nonabelian chief kernels
with their nontrivial permutation action and all
prescribed peripheral base conjugators.  It is not
asserted cofinal among the finite quotients relevant
to D.1 or D.2.  For example, the five-column cover
underlying D.27 has boundary counts $1,5,1$ and genus
zero; its relation has no handle commutators for the
count in \textup{(D.28.8)}.  The explicit maps of
D.27 and D.28 specify the respective finite lifting
conclusions without resolving the remaining cofinal
comparison.
